\chapter{Cơ Sở Lý Thuyết}
\label{chap:theory}

\section{Knowledge Tracing}
\label{sec:kt}

Knowledge Tracing (KT) là bài toán theo dõi và dự đoán trạng thái kiến thức của người học qua thời gian, dựa trên lịch sử tương tác với các bài tập. Cho một chuỗi tương tác $\{(q_1, r_1), (q_2, r_2), \ldots, (q_t, r_t)\}$ trong đó $q_t$ là câu hỏi tại thời điểm $t$ và $r_t \in \{0, 1\}$ là kết quả (đúng/sai), mục tiêu của KT là dự đoán xác suất trả lời đúng câu hỏi tiếp theo $q_{t+1}$:
\begin{equation}
  p(r_{t+1} = 1 \mid q_{t+1}, q_1, r_1, \ldots, q_t, r_t)
\end{equation}

Trong hệ thống \system{}, bài toán KT được áp dụng để theo dõi năng lực của ứng viên trên 17 nhóm kỹ năng (Knowledge Components -- KC) thuộc ba vai trò nghề nghiệp IT: Data Analyst (DA), Data Scientist (DS) và Data Engineer (DE).

\subsection{Bayesian Knowledge Tracing (BKT)}
\label{subsec:bkt}

BKT \cite{corbett1994knowledge} mô hình hóa từng KC như một chuỗi Markov ẩn (Hidden Markov Model -- HMM) với hai trạng thái: \textit{chưa thành thạo} ($L=0$) và \textit{thành thạo} ($L=1$). Mô hình được đặc trưng bởi bốn tham số cho mỗi KC:

\begin{itemize}
  \item $P(L_0)$: xác suất ban đầu người học đã thành thạo KC.
  \item $P(T)$: xác suất chuyển từ chưa thành thạo sang thành thạo sau mỗi lần thực hành.
  \item $P(S)$: xác suất trả lời sai dù đã thành thạo (\textit{slip}).
  \item $P(G)$: xác suất trả lời đúng dù chưa thành thạo (\textit{guess}).
\end{itemize}

Xác suất trả lời đúng tại bước $t$ được tính:
\begin{equation}
  P(r_t = 1) = P(L_t)(1 - P(S)) + (1 - P(L_t))P(G)
\end{equation}

Sau mỗi quan sát, trạng thái thành thạo được cập nhật theo Bayes:
\begin{equation}
  P(L_t \mid r_t = 1) = \frac{P(L_t)(1 - P(S))}{P(r_t = 1)}
\end{equation}

\begin{equation}
  P(L_{t+1}) = P(L_t \mid r_t) + (1 - P(L_t \mid r_t)) \cdot P(T)
\end{equation}

Tham số của BKT được ước lượng bằng thuật toán Expectation-Maximization (EM) hoặc tối ưu hóa trực tiếp L-BFGS-B trên log-likelihood. Ưu điểm của BKT là khả năng diễn giải cao và hiệu quả trên dữ liệu nhỏ; hạn chế là giả định độc lập giữa các KC và không mô hình hóa được tương tác bài tập phức tạp.

\subsection{Deep Knowledge Tracing (DKT)}
\label{subsec:dkt}

DKT \cite{piech2015deep} thay thế mô hình xác suất của BKT bằng mạng nơ-ron hồi quy (LSTM), cho phép học được biểu diễn ẩn phong phú hơn cho trạng thái kiến thức của người học. Đầu vào tại mỗi bước $t$ là vector one-hot mã hóa cặp $(q_t, r_t)$, có chiều $2 \cdot |\mathcal{Q}|$ với $|\mathcal{Q}|$ là số lượng câu hỏi (hoặc KC).

\begin{equation}
  \mathbf{h}_t = \text{LSTM}(\mathbf{x}_t, \mathbf{h}_{t-1})
\end{equation}
\begin{equation}
  \hat{y}_{t+1} = \sigma(\mathbf{W} \cdot \mathbf{h}_t + \mathbf{b})
\end{equation}

trong đó $\hat{y}_{t+1}$ là xác suất trả lời đúng câu hỏi tiếp theo, $\sigma$ là hàm sigmoid.

Trong \system{}, DKT được mở rộng với \textbf{quality label}: thay vì nhãn nhị phân $r_t \in \{0,1\}$, mô hình sử dụng điểm chất lượng liên tục $r_t \in [0, 1]$ phản ánh mức độ thành thạo (điểm tự đánh giá sau phỏng vấn). Hàm mất mát chuyển từ Binary Cross-Entropy sang Mean Squared Error:
\begin{equation}
  \mathcal{L}_{\text{quality}} = \frac{1}{T} \sum_{t=1}^{T} (r_t - \hat{y}_t)^2
\end{equation}

\subsection{Self-Attentive Knowledge Tracing (SAKT)}
\label{subsec:sakt}

SAKT \cite{pandey2019self} thay thế kiến trúc hồi quy của DKT bằng cơ chế self-attention, cho phép mô hình học được mối liên hệ trực tiếp giữa các KC liên quan trong lịch sử tương tác.

Cho chuỗi tương tác $\{(q_1, r_1), \ldots, (q_{t-1}, r_{t-1})\}$, SAKT xây dựng:
\begin{itemize}
  \item \textbf{Exercise embedding}: $\mathbf{e}_{q_t} \in \mathbb{R}^d$ cho câu hỏi cần dự đoán.
  \item \textbf{Interaction embedding}: $\mathbf{m}_s = \mathbf{e}_{q_s} + \mathbf{r}_s$ cho mỗi bước lịch sử.
\end{itemize}

Attention score giữa câu hỏi $q_t$ và bước lịch sử $s$ được tính:
\begin{equation}
  a_{t,s} = \frac{(\mathbf{e}_{q_t} \mathbf{W}^Q)(\mathbf{m}_s \mathbf{W}^K)^\top}{\sqrt{d}}
\end{equation}

\begin{equation}
  \hat{y}_t = \sigma\left(\mathbf{W}_o \cdot \sum_s \text{softmax}(a_{t,s}) \cdot \mathbf{m}_s \mathbf{W}^V\right)
\end{equation}

Ưu điểm của SAKT là khả năng song song hóa (không sequential như LSTM) và diễn giải được attention weight để hiểu KC nào ảnh hưởng đến dự đoán.

\section{Item Response Theory (IRT)}
\label{sec:irt}

Item Response Theory \cite{embretson2000item} là khung lý thuyết mô hình hóa mối quan hệ giữa năng lực người học $\theta$ và xác suất trả lời đúng câu hỏi $i$. Mô hình 3PL (3-Parameter Logistic) được dùng trong simulation:

\begin{equation}
  P(r = 1 \mid \theta, a_i, b_i, c_i) = c_i + (1 - c_i) \cdot \frac{1}{1 + e^{-a_i(\theta - b_i)}}
\end{equation}

trong đó $a_i$ là độ phân biệt (discrimination), $b_i$ là độ khó (difficulty), $c_i$ là xác suất đoán ngẫu nhiên (guessing). Trong \system{}, tham số IRT được ánh xạ từ cấu hình độ khó EASY/MEDIUM/HARD với $b \in \{0.30, 0.55, 0.80\}$.

\section{Recommendation System cho Interview Practice}
\label{sec:recsys}

Hệ thống gợi ý trong \system{} kết hợp KT và Collaborative Filtering để đề xuất câu hỏi phù hợp với trình độ hiện tại của ứng viên. Cụ thể, sau mỗi phiên luyện tập, vector năng lực $\bm{\theta}_u \in \mathbb{R}^{17}$ của người dùng $u$ (17 KC) được cập nhật bởi KT model, sau đó được dùng để xếp hạng câu hỏi theo chiến lược \textit{zone of proximal development}: ưu tiên câu hỏi thuộc KC yếu và độ khó phù hợp.

\section{Cập Nhật Năng Lực Thời Gian Thực: EMA và Blend KT--EMA}
\label{sec:ema}

Mô hình KT cho dự đoán tốt khi đã có đủ lịch sử tương tác, nhưng kém ổn định ở giai đoạn đầu (\textit{cold-start}) khi người dùng mới chỉ trả lời vài câu. Để giải quyết, \system{} kết hợp KT với \textbf{Exponential Moving Average (EMA)} -- một bộ ước lượng trực tuyến đơn giản, ổn định và không cần lịch sử dài.

Với mỗi KC, sau khi quan sát điểm số $s_t \in [0,1]$ tại bước $t$, giá trị EMA được cập nhật:
\begin{equation}
  \theta^{\text{EMA}}_t = \alpha \, s_t + (1 - \alpha)\,\theta^{\text{EMA}}_{t-1},
  \qquad \alpha = 0.35 .
\end{equation}

EMA luôn chạy ngay từ tương tác đầu tiên, cung cấp ước lượng năng lực thô. Khi lịch sử đạt tối thiểu $H_{\min} = 3$ tương tác, dự đoán KT được kích hoạt và \textbf{blend} với EMA theo trọng số cố định:
\begin{equation}
  \hat{\theta}_t = \lambda \, \theta^{\text{KT}}_t + (1 - \lambda)\,\theta^{\text{EMA}}_t,
  \qquad \lambda = 0.70 .
  \label{eq:blend}
\end{equation}

Công thức~\eqref{eq:blend} là một trong các đóng góp chính của khóa luận. Lựa chọn $\lambda = 0.7$ cho KT trọng số trội (vì KT nắm bắt được động lực chuỗi và quan hệ giữa các KC), trong khi giữ $0.3$ cho EMA nhằm \textit{ổn định hóa} dự đoán, chống dao động mạnh khi KT gặp chuỗi nhiễu. Ngưỡng $H_{\min}=3$ đảm bảo KT chỉ tham gia khi đã có đủ ngữ cảnh, còn trước đó hệ thống dùng EMA thuần. Hằng số $\lambda$ và $H_{\min}$ được giữ \textbf{nhất quán} giữa backend Python và frontend để bảo đảm con số hiển thị đồng bộ trên toàn hệ thống (xem Chương~\ref{chap:system}).

\section{Các Công Trình Liên Quan}
\label{sec:related}

\textbf{Knowledge Tracing.} Từ BKT \cite{corbett1994knowledge}, các mô hình KT đã phát triển qua DKT \cite{piech2015deep}, DKVMN \cite{zhang2017dynamic}, AKT \cite{ghosh2020context}, và gần đây là các mô hình Transformer-based như SAINT \cite{choi2020towards}. SAKT \cite{pandey2019self} là mô hình attention đầu tiên đơn giản và hiệu quả trong bài toán này.

\textbf{Adaptive Learning Systems.} Knewton \cite{ferreira2017knewton} và các hệ thống ITS (Intelligent Tutoring System) như ASSISTments \cite{heffernan2014assistments} sử dụng KT để cá nhân hóa lộ trình học. Tuy nhiên, các hệ thống này tập trung vào môi trường học thuật, chưa có ứng dụng cụ thể cho phỏng vấn IT.

\textbf{Interview Preparation Systems.} Các nền tảng như LeetCode, HackerRank cung cấp bài tập lập trình nhưng thiếu cơ chế theo dõi kiến thức có cấu trúc. \system{} khác biệt ở chỗ tích hợp KT với question bank đa dạng (lý thuyết, thực hành, coding) cho ba vai trò DA/DS/DE, đồng thời sử dụng LLM để chấm điểm tự động.
